% ===================================================================
%  BAGAN No. 12 — Stasiun. --- Jalan kereta api. --- Kapal.
%
%  Traduction, faite en 2026, en indonesien d'aujourd'hui, sur l'ido
%  de texto/io. Regles de la colonne : voir 10-bagan-01.tex. Une seule
%  numerotation. Le « Noto » d'ouverture appartient a l'apparat et n'a
%  pas de cle propre, comme du cote ido ; la note du feuillet 85 en a
%  une, t12-noto-1, et le bloc t12-14-1 se poursuit apres elle sous la
%  MEME cle.
%
%  DEUX CHEMINS DE FER COLONIAUX, DEUX LANGUES SOURCES, ET RIEN
%  D'AUTRE NE DIFFERE. La colonne ourdoue avait montre que le chemin
%  de fer indien laisse sept emprunts anglais — انجن, پلیٹ فارم, ٹکٹ,
%  گارڈ, سگنل, بوائلر, پسٹن — et que ce sont exactement les mots qu'un
%  voyageur lit ECRITS sur place. L'indonesien donne la meme liste, au
%  mot pres, mais en neerlandais : lokomotif, peron, karcis (de
%  kaartjes), kondektur, sinyal, ketel, piston, plus wesel l'aiguille
%  et sepur la voie. Les Pays-Bas ouvrent la premiere ligne de Java en
%  1867, la Grande-Bretagne celle du Bengale en 1854 : deux
%  administrations, deux lexiques, et le MEME decoupage. Ce n'est donc
%  pas la langue qui decide de ce qui s'emprunte, c'est l'ETIQUETTE
%  posee sur le quai.
%
%  kereta api, LE TRAIN : « LA VOITURE DE FEU ». La ou l'ourdou dit
%  ریل گاڑی en empruntant « rail », l'indonesien compose sur deux mots
%  malais et rejoint sans le savoir le chinois 火车 et le japonais
%  汽車, qui disent la meme chose. Le train est l'un des rares objets
%  industriels que plusieurs langues d'Asie ont prefere DECRIRE plutot
%  qu'emprunter — sans doute parce qu'il se voyait de loin avant qu'on
%  ait a le nommer.
%
%  ET kuli (16) EST LE MEME MOT QU'AU TABLEAU 12 OURDOU. La colonne
%  ourdoue notait que قلی est ourdou et hindi, que l'anglais en a fait
%  « coolie » et l'a exporte dans le monde entier. Le voici arrive a
%  Java, dans une gare neerlandaise, ecrit kuli. Le mot a fait le tour
%  par l'anglais et il est revenu en Asie : c'est le premier terme du
%  releve qu'on retrouve dans DEUX colonnes non europeennes sans
%  qu'aucune des deux ne l'ait pris a l'autre.
% ===================================================================

%%K t12-apar-1 apar
\VUsaut{4.0mm}
\VUtitre{15.0pt}{60}{BAGAN No. 12}
\VUsaut{2.4mm}
\VUpk{11.4pt}{40}{Stasiun. --- Jalan kereta api. --- Kapal.}
\VUsaut{3.4mm}
\textsc{Catatan.} --- \textit{Daftar jadwal yang dipakai di setiap
negeri berbeda-beda ; karena itu kami memakai sebagai contoh
jam-jam} \VUgras{daftar Prancis.}

%%K t12-01-1 p
1. --- Saya baru saja melakukan perjalanan melintasi Jerman. Sebelum
berangkat, saya membeli sebuah \VUgras{buku jadwal}
\textsuperscript{(15)} untuk mengetahui jam keberangkatan kereta.
Setelah memastikan bahwa kereta yang paling cocok akan berangkat dari
Paris pukul sembilan malam dan tiba di Strasburg pukul satu lewat tiga
belas menit, saya menyiapkan perjalanan saya, dan keesokan harinya
saya minta diantar ke stasiun.

%%K t12-02-1 p
2. --- Ketika saya masuk ke ruang keberangkatan yang besar, di
belakang seorang \VUgras{pastor} \textsuperscript{(2)} desa tua yang
menjinjing \VUgras{tas perjalanannya} \textsuperscript{(3)}, banyak
\VUgras{penumpang} \textsuperscript{(1)} sudah tiba.

%%K t12-02-2 p
Karena saya hendak mengirim sebuah \VUgras{telegram}
\textsuperscript{(5)}, saya minta seorang \VUgras{pengawas}
\textsuperscript{(6)} menunjukkan kepada saya \VUgras{kantor telegraf}
\textsuperscript{(4)}.

%%K t12-03-1 p
3. --- Sebelum masuk ke \VUgras{ruang tunggu} \textsuperscript{(7)},
saya pergi membeli \VUgras{karcis} \textsuperscript{(9)} pulang pergi.

%%K t12-04-1 p
4. --- Saya tidak lama menunggu di \VUgras{loket}
\textsuperscript{(8)}, sebab hanya sedikit orang di sana. Seorang
\VUgras{prajurit} \textsuperscript{(10)}, yang sedang berangkat cuti,
berbaik hati memberikan gilirannya kepada saya, sehingga saya sempat
menanyakan beberapa keterangan kepada \VUgras{kepala stasiun}
\textsuperscript{(13)}, yang sedang berbincang dengan
\VUgras{komisaris} \textsuperscript{(12)} pengawasan dan dengan
\VUgras{polisi militer} \textsuperscript{(11)} yang bertugas.

%%K t12-tit-1 sub
\VUpk{10.2pt}{40}{Pendaftaran bagasi.}
\VUsaut{3.0mm}

%%K t12-05-1 p
5. --- Sesudah membeli koran di \VUgras{kios buku}
\textsuperscript{(14)}, saya minta \VUgras{bagasi} saya
\textsuperscript{(17)} ditimbang — bagasi yang baru saja dibawa
seorang \VUgras{kuli} \textsuperscript{(16)} dengan \VUgras{gerobak
barang} \textsuperscript{(19)}. \VUgras{Juru tulis}
\textsuperscript{(22)} yang bertugas pada \VUgras{timbangan}
\textsuperscript{(21)} memberitahu saya bahwa koper saya beratnya tiga
puluh lima kilogram ; itu berarti \VUgras{kelebihan berat} lima
kilogram, yang untuknya saya harus membayar tambahan di \VUgras{loket
bagasi} \textsuperscript{(23)}.

%%K t12-06-1 p
6. --- Karena saya datang terlalu awal, saya sempat mengamati lalu
lalang para penumpang di stasiun. Sebagian menitipkan atau mengambil
kembali barang kecil mereka di \VUgras{tempat penitipan}
\textsuperscript{(25)} ; yang lain melihat \VUgras{daftar jadwal}
\textsuperscript{(26)}. Penumpang yang baru tiba bergegas ke
\VUgras{pintu keluar} \textsuperscript{(32)} — dengan \VUgras{koper}
\textsuperscript{(24)} di tangan. Beberapa orang menunggu di
\VUgras{tempat pengambilan bagasi} \textsuperscript{(27)} dan
menunjukkan \VUgras{resi} \textsuperscript{(29)} mereka untuk mengambil
kembali koper, \VUgras{peti} \textsuperscript{(28)} atau
\VUgras{keranjang} \textsuperscript{(30)} mereka.

%%K t12-07-1 p
7. --- \VUgras{Para petugas bea} \textsuperscript{(33)} memeriksa
bungkusan dengan cermat untuk memastikan apakah ada sesuatu yang harus
dilaporkan.

%%K t12-08-1 p
8. --- Beberapa penumpang pergi ke \VUgras{restorasi}
\textsuperscript{(34)}, tempat seorang \VUgras{pelayan}
\textsuperscript{(35)} bersemangat melayani mereka. Mereka harus
bergegas, sebab \VUgras{kereta api} \textsuperscript{(36)}, yang
merupakan kereta ekspres, tidak akan lama berhenti di stasiun.

%%K t12-09-1 p
9. --- Kemudian saya pergi ke peron keberangkatan dan mencari
\VUgras{gerbong} \textsuperscript{(37)} langsung supaya tidak terpaksa
berganti gerbong selama perjalanan. Saya naik ke sebuah kereta
berkoridor yang mempunyai satu \VUgras{kompartemen}
\textsuperscript{(38)} untuk perokok dan satu kompartemen yang
disediakan untuk « nyonya-nyonya yang berpergian sendiri ».

%%K t12-tit-2 sub
\VUpk{10.2pt}{40}{Berangkat pada malam hari.}
\VUsaut{3.0mm}

%%K t12-10-1 p
10. --- Sesudah menaiki \VUgras{injakan} \textsuperscript{(40)},
menutup \VUgras{pintu} \textsuperscript{(39)}, menggulung \VUgras{tirai}
\textsuperscript{(41)} dan meletakkan barang kecil saya di
\VUgras{jaring bagasi} \textsuperscript{(43)}, saya duduk di
\VUgras{bangku} \textsuperscript{(42)}. Melihat \VUgras{tukang lampu}
\textsuperscript{(44)} menyalakan lampu, saya berpikir bahwa saya harus
bermalam di gerbong, dan saya memanggil petugas yang mengurus
penyewaan \VUgras{bantal kepala} \textsuperscript{(45)} untuk meminta
sebuah.

%%K t12-11-1 p
11. --- \VUgras{Kondektur} datang memeriksa karcis. Saya bertanya
kepadanya mengapa kami belum juga berangkat, padahal waktunya sudah
tiba. Ia menjawab bahwa \VUgras{kepala kereta} \textsuperscript{(46)}
sedang sibuk memasukkan sepeda ke \VUgras{gerbong barang}
\textsuperscript{(47)}. Lagi pula, belum semua \VUgras{penghangat kaki}
\textsuperscript{(48)} diganti.

%%K t12-12-1 p
12. --- Tetapi kami segera akan berangkat, sebab \VUgras{juru api}
\textsuperscript{(50)}, di atas \VUgras{gerbong batu baranya}
\textsuperscript{(49)}, mengambil batu bara untuk membangkitkan api
\VUgras{lokomotif} \textsuperscript{(54)}, dan \VUgras{masinis}
\textsuperscript{(51)} hanya menunggu isyarat \VUgras{wakil kepala
stasiun} \textsuperscript{(52)} yang berdiri di \VUgras{peron}
\textsuperscript{(53)}.

%%K t12-13-1 p
13. --- \VUgras{Ketel} \textsuperscript{(56)} lokomotif menggeram
tertahan ; \VUgras{cerobong} \textsuperscript{(55)} memuntahkan arus
asap bercampur \VUgras{uap} \textsuperscript{(60)}. Sebuah
\VUgras{peluit} \textsuperscript{(59)} yang melengking berbunyi nyaring
dan \VUgras{piston} \textsuperscript{(58)} mulai bergerak, mendorong
\VUgras{roda} \textsuperscript{(57)} mesin yang berat itu.

%%K t12-tit-3 sub
\VUsaut{3.0mm}
\VUpk{10.2pt}{40}{Berangkat!}
\VUsaut{3.0mm}

%%K t12-14-1 p
14. --- Kecepatan kereta yang berlari di atas \VUgras{rel}
\textsuperscript{(64)} makin bertambah ; \VUgras{peron-peron}
\textsuperscript{(65)} lenyap ; kami melewati \VUgras{jamban}
\textsuperscript{(66)}, dan juga rangkaian gerbong yang diparkir di
\VUgras{jalur} \textsuperscript{(63)} yang lain : \VUgras{gerbong-
gerbong tidur} \textsuperscript{(67)}, \VUgras{gerbong makan}
\textsuperscript{(68)}, \VUgras{gerbong pos} \textsuperscript{(69)}.

%%K t12-noto-1 noto
\VUnotes{143.2mm}{%
\textit{(*) Catatan.} --- \textit{Sudah barang tentu gambaran
perjalanan ini mengikuti perbatasan sebelum perang.}%
}

%%K t12-15-1 p
15. --- Kemudian \VUgras{stasiun barang} \textsuperscript{(71)} lewat
di depan mata kami. Di bawah hanggar, para petugas memuat
\VUgras{barang} \textsuperscript{(72)} yang dikirim dengan kereta
lambat. \VUgras{Regu-regu langsir} \textsuperscript{(76)} di dekat
\VUgras{menara air} \textsuperscript{(73)} melangsir \VUgras{gerbong-
gerbong ternak} \textsuperscript{(75)} dan \VUgras{gerbong-gerbong
datar} \textsuperscript{(79)} untuk membentuk sebuah \VUgras{kereta
barang} \textsuperscript{(80)}.

%%K t12-16-1 p
16. --- Kami melewati \VUgras{wesel} \textsuperscript{(78)} yang
terakhir, tempat tukang wesel berdiri ; kami melewati \VUgras{sinyal}
\textsuperscript{(86)} dan stasiun itu lenyap di kejauhan.

%%K t12-17-1 p
17. --- Segera kami melintasi sebuah \VUgras{jembatan besi}
\textsuperscript{(74)} yang menyeberangi \VUgras{sungai}
\textsuperscript{(81)}. Sesudah melintasi jembatan itu, kami menyusuri
sungai, tempat \VUgras{kapal-kapal tunda} \textsuperscript{(82)} yang
perkasa menarik \VUgras{tongkang-tongkang} \textsuperscript{(83)} yang
berat. Kami juga melihat sebuah \VUgras{kapal penumpang}
\textsuperscript{(84)} ketika ia merapat ke sebuah \VUgras{ponton}
\textsuperscript{(85)}.

%%K t12-18-1 p
18. --- Sesudah melewati beberapa stasiun dengan sangat cepat, kami
melintasi beberapa \VUgras{terowongan} \textsuperscript{(87)} dan
meninggalkan di belakang kami \VUgras{rumah-rumah kecil}
\textsuperscript{(88)} \VUgras{penjaga palang} yang tak terhitung
banyaknya. Dibuai oleh goyangan kereta, saya tertidur nyenyak.

%%K t12-tit-4 sub
\VUsaut{3.0mm}
\VUpk{10.2pt}{40}{Stasiun Deutsch-Avricourt.}
\VUsaut{3.0mm}

%%K t12-19-1 p
19. --- Saya tiba-tiba dibangunkan oleh suara seorang petugas yang
berteriak : « Deutsch-Avricourt! \textsuperscript{(*)} Semua turun! »
Terjadilah pemeriksaan bea cukai dan segala urusan perbatasan yang
membosankan. Saya harus menyesuaikan arloji saku saya dengan
\VUgras{jam besar} \textsuperscript{(89)} stasiun, yang lima puluh lima
menit lebih cepat ; baru saja bangun, saya sulit membedakan
\VUgras{jarum besar} \textsuperscript{(90)} dari \VUgras{jarum kecil}
\textsuperscript{(91)} ; dan \VUgras{piringan angkanya}
\textsuperscript{(92)} sendiri sama sekali kabur karena kabut pagi.

%%K t12-20-1 p
20. --- Sementara para petugas bea mengadakan pemeriksaan, saya sambil
menguap mengeja \VUgras{poster-poster} \textsuperscript{(93)} yang
menghiasi dinding stasiun. Untuk mengisi waktu saya sarapan, melihat
peta \VUgras{jaringan} \textsuperscript{(94)} kereta api Jerman untuk
mengetahui berapa jauh lagi Strasburg dari saya, dan saya masuk
kembali ke gerbong saya, dikuatkan oleh harapan segera mencapai tujuan
perjalanan saya.

\VUsaut{6.0mm}
\VUfilet{22mm}
